% !PRINTED 303
\weberpage{303}

But we can also give up the assumption that none of the functions (1)
vanishes for $\alpha$ and $\beta$, provided that $\alpha$ and $\beta$ do
not occur among the roots of $f(x)$.

For then our theorem is applicable to two values, of which the first is a
little greater than $\alpha$, the other a little smaller than $\beta$, and
which enclose the same roots as $\alpha$ and $\beta$; that is, if we
denote by $\mathfrak{z}$ the number of these roots, we have
\[
V(\alpha + 0) - V(\beta - 0) = \mathfrak{z} + 2h,
\]
and since
\[
V(\alpha + 0) = V(\alpha), \quad V(\beta - 0) = V(\beta) + 2h',
\]
it follows that
\[
V(\alpha) - V(\beta) = \mathfrak{z} + 2(h + h'),
\]
which is again the statement of our theorem.

But if $\alpha$ or $\beta$ itself belongs to the roots of $f(x)$, then our
theorem is correct only \gesperrt{when these limiting values are not counted as
belonging to the interval}.

The \gesperrt{Budan-Fourier} theorem does not, it is true, give a
completely certain decision about the number of roots in an interval, as
the \gesperrt{Sturm} theorem does; but it can nevertheless in many cases
replace the \gesperrt{Sturm} theorem; for if one has restricted the
interval $(\alpha, \beta)$ so far that no change of sign, or only one, is
lost from $\alpha$ to $\beta$, then it follows with certainty that in the
first case no root, and in the second case one and only one root, lies in
the interval.

The separation of the roots can therefore be completely given by the Budan
sequence only when, in the passage through a value $\xi$, several changes
of sign are not lost simultaneously. This behaviour cannot be recognised
with certainty, however small the interval $(\alpha, \beta)$ may be
chosen, and so, if after a suitable narrowing of the interval the
separation has not succeeded, one will after all have to resort to another
procedure, in the last instance to the \gesperrt{Sturm} theorem.
